%This is a template for producing OASIcs conference volumes.
%The usage of this file together with oasicsmaster-v2019.cls should be
%straightforward. There is no separate documentation.

\documentclass[a4paper,UKenglish]{oasicsmaster-v2021}
%for A4 paper format use option "a4paper", for US-letter use option "letterpaper"
%for british hyphenation rules use option "UKenglish", for american hyphenation rules use option "USenglish"

%\graphicspath{{./graphics/}}%helpful if your graphic files are in another directory

\bibliographystyle{plainurl}% the mandatory bibstyle

\editor{John Q. Public}{Dummy University Computing Laboratory, [Address], Country \and Second affiliation, Country}{johnqpublic@dummyuni.org}{https://orcid.org/0000-0002-1825-0097}%TODO mandatory, please use full name; only 1 editor per \editor macro; first two parameters are mandatory, other parameters can be empty.

\editor{Joan R. Access}{School of Computer Science, University City, Country2}{beditor@uni-city.edu}{}

%TODO please enter name of conference without abbreviate title 
\EventTitle{42nd Conference on Very Important Topics}

%TODO please enter macros as provided by OASIcs office%%%%%%%%%%%%%%%%%%%%%%%%%%%%%%%%%%
\EventEditors{John Q. Open and Joan R. Access}
\EventNoEds{2}
\EventLongTitle{42nd Conference on Very Important Topics (CVIT 2016)}
\EventShortTitle{CVIT 2016}
\EventAcronym{CVIT}
\EventYear{2016}
\EventDate{December 24--27, 2016}
\EventLocation{Little Whinging, United~Kingdom}
\EventLogo{}
\SeriesVolume{42}
\ArticleNo{0}
%%%%%%%%%%%%%%%%%%%%%%%%%%%%%%%%%%%%%%%%%%%%%%%%%%%%%%

\ccsdesc[100]{\textcolor{red}{Replace ccsdesc macro with valid one}}%TODO mandatory: Please choose ACM 2012 classifications from https://dl.acm.org/ccs/ccs_flat.cfm . 

\ISBN{424-2-424242-424-2}%TODO mandatory: Please enter ISBN as provided by OASIcs office

%\DatePublished{June 6, 2018}%leave empty; will be entered by OASIcs office before publication

%\title{\huge 1st Conference on Very Important Topics} %use only to override default title

%\subtitle{CVIT 2016, December 24--27, 2016, Little Whinging, United Kingdom} %use only to override default subtitle

%\titlepagebottomline{OASIcs -- Vol.~\printSeriesVolume{} -- \printEventShortTitle \qquad \qquad \qquad www.dagstuhl.de/oasics} %override text line in the yellow box on the titlepage

%\serieslogo{} %please provide filename (without suffix)

%\nolinenumbers %uncomment to disable line numbering

\begin{document}

\frontmatter

%%
%% PAGE 1: Cover page
%%%

\maketitle

%%
%% PAGE 2: Bibliographic data (editors, ACM classification, ISBN, license, DOI, ...)
%% 

\begin{publicationinfo}%for page ii, please fill as required
\sffamily

\emph{Editors} 
\printEditorLong


\bigskip
\bigskip
\emph{ACM Classification 2012}\\
\printSubjclass

\bigskip
\bigskip

{\Large\bfseries\sffamily \href{https://www.dagstuhl.de/dagpub/\printISBN}{ISBN \printISBN}}

\bigskip
\bigskip

\emph{Published online and open access by}\newline
Schloss Dagstuhl -- Leibniz-Zentrum f\"ur Informatik GmbH, Dagstuhl Publishing, Saarbr\"ucken/Wadern, Germany. Online available at \href{https://www.dagstuhl.de/dagpub/\printISBN}{https://www.dagstuhl.de/dagpub/\printISBN}.

\bigskip
\emph{Publication date}\newline
\printDatePublished

\bigskip
\bigskip

\emph{Bibliographic information published by the Deutsche Nationalbibliothek}\newline
The Deutsche Nationalbibliothek lists this publication in the Deutsche Nationalbibliografie; detailed bibliographic data are available in the Internet at \href{https://portal.dnb.de}{https://portal.dnb.de}. 

\bigskip

\emph{License}\newline
This work is licensed under a Creative Commons Attribution 4.0 International license (CC-BY~4.0):\\ \href{https://creativecommons.org/licenses/by/4.0/legalcode}{https://creativecommons.org/licenses/by/4.0/legalcode}.\\
In brief, this license authorizes each and everybody to share (to copy, distribute and transmit) the work under the following conditions, without impairing or restricting the authors' moral rights:
\marginpar{\hspace*{0.2\marginparwidth}\includegraphics[width=0.75\marginparwidth]{cc-by.pdf}}
\begin{itemize}
\item Attribution: The work must be attributed to its authors.
\end{itemize}

\smallskip

The copyright is retained by the corresponding authors.

\bigskip
\bigskip
\bigskip
\bigskip

Digital Object Identifier: \href{https://doi.org/\printDOI}{\printDOI}

\vfill
\textbf{\href{https://www.dagstuhl.de/dagpub/\printISBN}{ISBN \printISBN}}\qquad \qquad \textbf{\href{https://www.dagstuhl.de/dagpub/1868-8969}{ISSN 1868-8969}}  \hfill \textbf{\href{https://www.dagstuhl.de/oasics}{https://www.dagstuhl.de/oasics}}

%%
%% PAGE 3: OASIcs series information
%% 

\newpage

\ \\
\bigskip
\bigskip
\bigskip

{\Large OASics -- OpenAccess Series in Informatics}

\bigskip

OASics is a series of high-quality conference proceedings across all fields in informatics. 
OASics volumes are published according to the principle of Open Access, i.e., they are available online and free of charge. 

\bigskip
\bigskip
\bigskip

\emph{Editorial Board}

%OASIcs Board members as of June 2018
\begin{itemize}
\item Daniel Cremers (TU M\"unchen, Germany)
\item Barbara Hammer (Universit\"at Bielefeld, Germany)
\item Marc Langheinrich (Universit\`a della Svizzera Italiana -- Lugano, Switzerland)
\item Dorothea Wagner (\emph{Editor-in-Chief}, Karlsruher Institut f\"{u}r Technologie, Germany)
 \end{itemize}


\bigskip
\bigskip
\bigskip

{\large\bfseries\sffamily \href{https://www.dagstuhl.de/dagpub/1868-8969}{ISSN 1868-8969}}

\bigskip
\bigskip
\bigskip

{\Large\bfseries\sffamily \href{https://www.dagstuhl.de/oasics}{https://www.dagstuhl.de/oasics}}

\vfill

%%
%% PAGE 4: (empty)
%% 

\newpage

\thispagestyle{empty}

\ \\

\end{publicationinfo}

%%
%% PAGE 5 and more: TOC etc.
%% 

%TODO please fill or comment out
\begin{dedication}
Insert dedication here.
\end{dedication}


\begin{contentslist}

%Please note, that the table of contents will be generated by the OASIcs office based on the order provided in the Dagstuhl Submission System. So you can keep this section as is!!!

%To generate the table of contents copy all the .vtc files
%of the contributions to your working directory.
%For every contribution type a line
%\inputtocentry{dummycontribution}
%where the argument of \inputtocentry is the name of
%the vtc file without suffix.

%Alternatively write e.g.
\contitem
\title{Preface}
\author{John Q. Open}
\page{0:vii}

%\part{} %use if volume is divided in parts
\part{Regular Papers}

\contitem
\title{Mmmmm $\ldots$ donuts}
\author{Homer J. Simpson}
\page{1:1--1:23}

\inputtocentry{oasics-v2021-sample-article}

\end{contentslist}


%TODO please fill or comment out
\chapter{Preface}

Lorem ipsum dolor sit amet, consectetur adipiscing elit. Ut mattis
elementum fermentum. Pellentesque habitant morbi tristique senectus et
netus et malesuada fames ac turpis egestas. Nulla sapien magna,
bibendum in dictum sed, egestas vel purus. Pellentesque id ornare
lacus. Pellentesque justo elit, sodales a fringilla vitae, gravida sed
elit. Etiam turpis eros, tincidunt sit amet tempor sed, gravida quis
eros. Mauris et nunc enim. Ut congue rhoncus odio vitae lacinia. Nunc
placerat est eu eros dignissim ac tristique nisi placerat. Fusce et
hendrerit justo. Nunc feugiat pulvinar nunc ac tincidunt. Donec eu
pharetra metus. Cras malesuada ante accumsan purus dignissim
euismod. Curabitur risus ante, aliquet ut suscipit eget, vulputate non
ligula. Nullam eleifend malesuada est, nec adipiscing sapien eleifend
eget. Vestibulum fringilla diam id felis sagittis aliquet. Maecenas
sed metus vel dui vulputate pretium et id lacus. Ut eget libero augue,
ut aliquet orci. Integer sed nunc id massa interdum imperdiet. Nunc ut
consequat eros.

Morbi hendrerit dapibus augue. Proin sed adipiscing ipsum. Ut
vulputate ultricies diam id dictum. Nunc pharetra imperdiet
sodales. Morbi convallis massa vitae justo adipiscing nec congue nunc
fringilla. Pellentesque eu rhoncus ligula. Nam eros neque, hendrerit a
rhoncus vel, molestie nec nulla. Ut iaculis vulputate mauris, non
scelerisque dolor fringilla sit amet. Pellentesque habitant morbi
tristique senectus et netus et malesuada fames ac turpis
egestas. Quisque vitae accumsan risus. Sed molestie dictum
venenatis. Sed odio justo, gravida et vulputate eu, congue et lorem. 


%optional
\begin{participants}
\chapter[Authors]{List of Authors}
%use \participant for every author, eg.:
\participant John Q. Public\\ 
Dummy University Computing Laboratory\\
Address, Country\\
johnqpublic@dummyuni.org

\end{participants} 


\end{document}
